\begin{table}[ht]
\footnotesize
\centering
\begin{threeparttable}
\caption{Final LSTM hyperparameters for the four forecasting feature sets.}
\label{tab:lstm_hparams}
\renewcommand{\arraystretch}{1.15}
\setlength{\tabcolsep}{6pt}
\begin{tabular}{lcccc}
\toprule
\textbf{Feature set} & \textbf{hidden} & \textbf{layers} & \textbf{dropout} & \textbf{lr} \\
\midrule
Price + volume & 128 & 2 & 0.205 & 0.004366 \\
Price + volume + sentiment & 128 & 2 & 0.138 & 0.002196 \\
Price + volume + technicals & 128 & 2 & 0.242 & 0.001519 \\
Full dataset & 64 & 1 & 0.249 & 0.003324 \\
\bottomrule
\end{tabular}
\begin{tablenotes}[flushleft]
\footnotesize
\item \textit{Note:} Selected by an Optuna study per feature set (50 trials with median pruning, each trial trained for 5 epochs on a reduced corpus of the first 10,000 training sequences, 5,000 for the price-and-volume set), scored by validation MAE on the inner-validation slice of the training period. The encoder is a unidirectional stacked LSTM with a linear head on the final hidden state. Batch size 256 and AdamW with the Huber loss on the standardised target are common to all runs.
\end{tablenotes}
\end{threeparttable}
\end{table}